\section{Genus-three real structures and the congruence-kernel fibre}

Write $\Gamma_g$ for the orientation-preserving mapping class group of
a closed oriented surface $S_g$, $\Gamma_g^{\pm}$ for its extension
by orientation reversal, $P_g=\widehat\Gamma_g$ for the full profinite
completion, and $Q_g=\check\Gamma_g$ for the closure of its action
in $\operatorname{Out}(\widehat{\pi_1(S_g)})$.  Superscripts $\pm$
on completions have the same orientation meaning.  The congruence
kernel is
\[
 K_c=\ker(q:P_3\longrightarrow Q_3).
\]
In particular $K_c$ is not the kernel of a homology representation.
We use $a=0$ for a dividing real structure and $a=1$ for a
nondividing one; $r$ is its number of ovals.

\statement{E.1}{Proved}{The compact real-type injection of LM6 C.13}
For $g\geq2$, the topological types of orientation-reversing
involutions of $S_g$ inject into conjugacy classes in
$\operatorname{Out}(\widehat{\pi_1(S_g)})$.  They consequently
inject into conjugacy classes in $Q_g^{\pm}$ and $P_g^{\pm}$.
For genus three the six types are
\[
 (r,a)=(0,1),(1,1),(2,1),(3,1),(2,0),(4,0).                 \tag{E.1.1}
\]

\begin{proof}
Put $\Pi=\widehat{\pi_1(S_g)}$.  The group $\Pi$ is centre-free.
An outer action of $C_2$ on $\Pi$ therefore determines its extension
uniquely: it is the pullback of
$\operatorname{Aut}(\Pi)\to\operatorname{Out}(\Pi)$, with its
kernel identified with $\Pi$ by inner automorphisms.  For a
geometric involution $\tau$ this is the completed orbifold group
of $S_g/\tau$, with its orientation subgroup $\Pi$.

We justify the finite-order comparison used for this particular
orbifold group.  Surface groups are good: after passing to a
finite cover to trivialize any given finite coefficient module,
degree-one classes are killed by the finite-index kernels of
their characters, and degree-two classes are killed by a further
cover of degree divisible by the exponent of the coefficient
group.  These comparisons, or the equivalent cohomological
effacement criterion, give the completion comparison in degrees
zero, one and two; higher degrees vanish.  The same argument
works for the orientation double of a nonorientable surface.
A finite extension preserves goodness, by the two
Hochschild--Serre spectral sequences and exactness of completion
for a finite-index subgroup.  The groups here have finite virtual
cohomological dimension two and finite cohomology in each degree
with finite coefficients.  Filtered unions of finite coefficient
modules give the comparison for all discrete continuous
$\mathbb F_2$-modules.  Thus all hypotheses of Theorem~A of
\cite{CEBoggiZalesskii2025} apply, at $p=2$, to this single
involution and its orbifold group.

The discrete order-two subgroups correspond exactly to the
ovals.  One way to see this is to choose an invariant hyperbolic
metric: a lift of a reflection has a geodesic axis, whose image
is an oval, and all lifts over the same oval are conjugate by
the surface group.  A lift over an oval has centralizer
$\mathbb Z\times C_2$, with the infinite cyclic factor generated
by a circuit of that oval.  For a fixed-point-free involution
the quotient is a closed nonorientable surface and its group
is torsion-free.  Theorem~A therefore identifies the completed
section classes with the $r$ ovals, and identifies their
centralizers in $\Pi$ with the full closures
$\widehat{\mathbb Z}$ of these cyclic groups.  Conjugacy by the
orbifold group gives the same section classes as conjugacy by
its orientation subgroup, since multiplying a conjugator by
the source reflection changes its orientation without changing
the conjugacy.

The image $w_i\in H_1(\Pi,\mathbb F_2)$ of a generator of the
$i$th oval centralizer is well defined: any other topological
generator differs by a profinite unit, whose reduction modulo
two is one.  The ovals divide the surface exactly when
\[
                         \sum_{i=1}^r w_i=0.                 \tag{E.1.2}
\]
Indeed a mod-two two-chain with that boundary has constant
coefficient on each component of the complement.  Its boundary
is all the ovals exactly when the two sides can be consistently
coloured with opposite coefficients.  This is precisely the
separating condition.  Equivalently, in the nondividing case
the connected complement would force one coefficient on both
sides of every oval.  If $r=0$, the structure is declared
nondividing, and this case is already detected by the absence
of sections.

Thus an isomorphism of the completed outer extensions preserves
$r$ and $a$.  Conversely those invariants classify the
topological involution: the quotient is a compact surface with
$r$ boundary components, orientable of genus
$(g+1-r)/2$ in the dividing case and nonorientable of genus
$g+1-r$ in the nondividing case.  Its orientation double in the
interior, completed along the mirror boundary, recovers the
surface with involution.  The classification of compact surfaces
therefore gives uniqueness.  In genus three these formulas give
exactly (E.1.1).  Conjugacy in either mapping class completion
would imply conjugacy of the outer actions, proving the two
remaining injections.
\end{proof}

\statement{E.2}{Proved}{Marked incidence and the simultaneous orientation in LM6 C.14}
The injection of E.1 holds for individually labelled real markings,
and for unordered markings allowing nonreal conjugate pairs.
A peripheral isomorphism of the completed real surface extensions
recovers the cyclic incidence of real marked points on each oval.
For dividing surfaces, the cyclic orders on labelled marked ovals
are identified only by their simultaneous reversal.

\begin{proof}
Let $U=S_g\setminus D$.  Filling the punctures is the quotient
of the completed geometric fundamental group by the closed
normal subgroup generated by its peripheral subgroups.  An
isomorphism preserving these conjugacy classes therefore induces
the compact comparison of E.1.  At a real puncture the local
extension is $\widehat{\mathbb Z}\rtimes_{-1}C_2$.  Its two
local section classes are distinguished by the parity of the
exponent in $c^t s$.  Their maps to global section classes say
which real intervals end at this puncture.  Consequently the
local-to-global incidence graph has real intervals as vertices,
real punctures as edges, and these two section classes as flags.
Unpunctured ovals give isolated compact components.  A punctured
oval gives its cyclic incidence graph.  Retaining flags is
essential even for a one-edge loop or a two-edge cycle.  This
description proves the claimed separate unoriented cyclic
orders directly from the extension.

We prove that the signs on marked ovals have their asserted
integral meaning; a procyclic generator cannot be replaced by
an arbitrary profinite unit.  Fix a marked oval, a base interval
next to a mark, and the corresponding pointed reflection $s$
after filling.  By E.1 its geometric centralizer is
$H=\overline{\langle a\rangle}\cong\widehat{\mathbb Z}$,
where $a$ makes one oriented circuit of the actual oval.
For an open normal $s$-invariant subgroup $N\leq\Pi$, the
subgroup $N\langle s\rangle$ defines the finite real cover with
the selected real base point.  Path lifting identifies the
fibre on the selected oval with
\[
 \langle a\rangle N/N=HN/N=H/(H\cap N).                    \tag{E.2.1}
\]
The inverse image of the marked oval graph is the corresponding
finite cyclic graph covering, including flags and the selected
sheet.  Elements of $H$ give deck transformations of this
cover because they commute with $s$ and normalize $N$.

The graph construction is functorial: both cusp flags map by
the local section maps, interval classes by composition of
sections, and all incidence maps commute.  A peripheral
extension isomorphism therefore gives compatible isomorphisms
of these graph coverings and of their deck actions.  The base
graph isomorphism has degree $\epsilon=\pm1$.  Lifting one
circuit at each finite level gives
\[
 b(a)\equiv(a')^\epsilon\pmod{H'\cap b(N)}.
\]
The sign is the same at every level because the graph coverings
lie over that same base graph map.  The subgroups $H\cap N$
are cofinal in the full profinite topology of $H$, so
$b(a)=(a')^\epsilon$ after the conjugations identifying the
chosen sections.

For a dividing surface orient all ovals as the boundary of
one half.  The integral relative homology sequence gives
\[
 \ker\left(\widehat{\mathbb Z}^{,r}
       \longrightarrow H_1(S_g,\widehat{\mathbb Z})\right)
       =\widehat{\mathbb Z}(1,\ldots,1).                     \tag{E.2.2}
\]
There are two complementary halves; their boundary vectors
are the displayed diagonal vector and its negative, which
proves that these are all relations.  An isomorphism sends
each oval generator to a unit multiple of the target oval
generator.  Preservation of (E.2.2) forces all these unit
multipliers to agree.  On every marked oval the multiplier is
the already determined sign $\epsilon$, so the marked signs
agree.  Nondividing quotient surfaces allow the independent
boundary reversals.  For unordered markings the nonreal
peripheral classes occur in two-element conjugation orbits;
these record the number of interior marks of the quotient.
The quotient surface with these interior marks and with the
recorded boundary cyclic incidence determines the topological
marked type.  Doubling its homeomorphisms gives the desired
classification upstairs.  This proves both marked injections.
\end{proof}

\statement{E.3}{Proved}{The exact nonabelian-cohomology fibre of LM6 C.15}
Let $1\to K\to P^{\pm}\xrightarrow{q}Q^{\pm}\to1$ be a
continuous exact sequence of profinite groups with compatible
orientation characters, with $K\leq P$ their positive kernel.
Fix an orientation-reversing involution $\tau\in P^{\pm}$ and
put $\alpha=\operatorname{Ad}(\tau)$ and $\bar\tau=q(\tau)$.
For $k\in K$ satisfying $k\alpha(k)=1$, identify cocycles by
$k\sim b k\alpha(b)^{-1}$, $b\in K$.  Then
\begin{align}
 q(g)\star[k]&=[gk\alpha(g)^{-1}],\notag\\
 \{\text{$P$-classes over }[\bar\tau]\}
      &\cong C_Q(\bar\tau)\backslash H^1(C_2,K),             \tag{E.3.1}\\
 1\longrightarrow K^{k\tau}\longrightarrow C_P(k\tau)
      &\xrightarrow{q}
       \operatorname{Stab}_{C_Q(\bar\tau)}([k])
       \longrightarrow1 .                                  \tag{E.3.2}
\end{align}
Here $g\in P$ is any lift of the indicated quotient-centralizer
element.  The action need not fix the class of the trivial cocycle.

\begin{proof}
The fibre over the element $\bar\tau$ consists precisely of
$k\tau$ with $k\alpha(k)=1$.  Conjugation by $b\in K$ gives
the indicated equivalence.  If $q(g)$ centralizes $\bar\tau$,
then $gk\alpha(g)^{-1}\in K$, and
\[
             g(k\tau)g^{-1}=gk\alpha(g)^{-1}\tau.
\]
This proves the cocycle condition and the action law.  Replacing
$g$ by $bg$ changes the cocycle by the stated equivalence, so
the action is independent of lifts.  An involution whose image
is conjugate to $\bar\tau$ can first be conjugated by a lift
in $P$ to lie over that element.  Two such representatives
are $P$-conjugate exactly when their $K$-classes are in the
same quotient-centralizer orbit.  This proves (E.3.1).

The kernel of the centralizer map is $K^{k\tau}$.  Its image
fixes $[k]$.  If a quotient-centralizer element fixes $[k]$,
choose its lift $g$.  There is $b\in K$ with
$g(k\tau)g^{-1}=b(k\tau)b^{-1}$, and $b^{-1}g$ is a
centralizing lift.  This proves (E.3.2), including surjectivity
onto the stated stabilizer.

For the last qualification take
$P^{\pm}=D_8=\langle r,t\mid r^4=t^2=1,trt=r^{-1}\rangle$,
$P=\langle r\rangle$, $K=\langle r^2\rangle$, and $\tau=t$.
The two classes in $H^1(C_2,K)$ are $[1],[r^2]$.  Although
all of $Q=P/K$ centralizes $\bar\tau$, its element $q(r)$
interchanges these classes because $r\alpha(r)^{-1}=r^2$.
Thus the orbit set has one element and $C_P(t)=K$ maps
trivially to $Q$.  The attached finite script checks these
group operations and orbit sizes.
\end{proof}

\statement{E.4}{Proved}{The hyperelliptic comparison of LM6 C.16 uses only the element theorem}
Fix a hyperelliptic involution $\iota\in\Gamma_3$, and put
$H^{\pm}=C_{\Gamma_3^{\pm}}(\iota)$ and $H=H^{\pm}\cap\Gamma_3$.
The natural map $\widehat{H^{\pm}}\to P_3^{\pm}$ is injective,
with image $\overline{H^{\pm}}$.  Each orientation-reversing
involution in this image is $\overline H$-conjugate to a
discrete one.  For such a discrete involution $\tau$,
\[
 C_{P_3}(\tau)\cap\overline H
       =\overline{C_H(\tau)}.                               \tag{E.4.1}
\]
Each of the six types in (E.1.1) has a representative in $H^{\pm}$.
No centralizer equality for a nonsolvable finite subgroup is
used in these assertions.

\begin{proof}
Theorem~7.1 of \cite{CEBoggiHyperelliptic2024} states the
congruence subgroup theorem with the topology induced from
the action on the completed surface group.  Thus the
composite $\widehat H\to P_3\to Q_3$ is injective.  The first
map is injective, and adjoining the orientation quotient of
order two proves injectivity for $H^{\pm}$ as well.

We check the group to which the torsion theorem is applied.
The branched double cover gives the Birman--Hilden exact
sequence
\[
 1\longrightarrow\langle\iota\rangle\longrightarrow H
       \longrightarrow\Gamma_{0,[8]}\longrightarrow1.
                                                               \tag{E.4.2}
\]
The kernel is the deck group.  All sphere mapping classes
lift because the double-cover character sends each of the
eight meridians to the nonidentity element of $C_2$.
This is the exact sequence (3) in the same primary source.
Intersect the inverse image of the pure sphere group with
the level-three homology kernel.  The resulting finite-index
subgroup misses $\iota$, which acts as $-I$ on
$H_1(S_3,\mathbb F_3)$, and embeds as a finite-index subgroup
of the pure sphere mapping class group.

Here is a direct goodness verification.  After fixing three
points at $0,1,\infty$, the forgetful configuration fibration
has finitely generated free fundamental group as kernel and
a section, obtained by placing the new real point at
$2+\sum_j|z_j|$.  A finite-index subgroup of a finitely
generated free kernel contains a characteristic finite-index
subgroup.  The action on its finite quotient has finite image;
the resulting finite semidirect-product quotient of the whole
group detects that kernel quotient.  Thus completion of the
split extension is exact.  Induction and Hochschild--Serre
comparison, starting from free groups, prove goodness of
the pure sphere group.  They also prove finite cohomology
with finite coefficients and finite virtual cohomological
dimension.  Shapiro's lemma passes these properties to
finite-index subgroups, and Hochschild--Serre passes them to
finite extensions.  They therefore hold for $H^{\pm}$.
Filtered unions again give the comparison for arbitrary
discrete continuous $\mathbb F_2$-modules.

Theorem~A of \cite{CEBoggiZalesskii2025}, at $p=2$, now
gives torsion conjugacy and centralizer equality inside
$\widehat{H^{\pm}}$ for the one element $\tau$.  A negative
conjugator is made positive by multiplying by the source
involution.  Intersecting its centralizer equality with the
open positive subgroup gives (E.4.1).

For existence of all types use the smooth projective models
$y^2=f(x)$ with
\[
\begin{array}{c|c}
 (r,a)&f(x)\\\hline
 (0,1)&-(x^8+1)\\
 (1,1)&(x^2-1)(x^6+1)\\
 (2,1)&(x^2-1)(x^2-4)(x^4+1)\\
 (3,1)&(x^2-1)(x^2-4)(x^2-9)(x^2+1)\\
 (2,0)&x^8+1\\
 (4,0)&\displaystyle\prod_{j=1}^4(x^2-j^2).
\end{array}                                                   \tag{E.4.3}
\]
All polynomials are square-free of degree eight.  If there
are $2k>0$ real roots and positive leading coefficient, the
$k$ positive intervals on $\mathbb P^1(\mathbb R)$ lift to
$k$ ovals.  For $k<4$ the inverse image of each open
half-plane is connected because it contains branch points;
a negative real interval joins these halves through nonreal
points of the curve.  The complement of the ovals is
therefore connected.  The four-oval case is dividing by the
quotient Euler characteristic in E.1.  The negative polynomial
has no real points, including at infinity.  For $x^8+1$, the
two signs give two real circles, as the infinity coordinate
$y/x^4$ preserves each sign.  All points above the real
projective line are real, and the two connected half-plane
preimages are the two complementary halves.  Conjugation
commutes with the hyperelliptic deck involution in all six
models.  Every genus-three hyperelliptic involution is
conjugate to $\iota$, by the topological double-cover
construction with eight branch points, which proves the
claim for the specified $H^{\pm}$.
\end{proof}

\statement{E.5}{Proved}{The five oval centralizers of LM6 C.17}
Let $\tau\in\Gamma_3^{\pm}$ be a real involution with $r>0$,
and put $C=C_{\Gamma_3}(\tau)$.  The restriction of $q$ is an
isomorphism
\[
 \overline C^{\,P_3}\xrightarrow{\sim}
 C_{Q_3}(\bar\tau)=\overline C^{\,Q_3},\qquad
 C_{P_3}(\tau)=K_c^\tau\rtimes\overline C^{\,P_3}.           \tag{E.5.1}
\]
For these five types the quotient-centralizer action of E.3
fixes the trivial cohomology class.  Formula (E.5.1) does not
compute $K_c^\tau$.

\begin{proof}
Let $A$ be the oval multicurve.  Its curves are essential and
pairwise nonisotopic.  A contractible oval would bound a disc
which is exchanged with its other side by the involution,
forcing genus zero.  Parallel ovals would similarly force
the complementary component on either side to be an annulus,
forcing genus one.  Neither can occur here.

First take $b\in C_{Q_3}(\bar\tau)$.  It permutes the section
classes of the real surface extension and their centralizers.
To recover actual curve generators, choose a real marked point
on an oval and use the pointed comparison
\[
 \check\Gamma_{3,1}^{\pm}\cong
 \operatorname{Aut}(\Pi)\mathbin{\times}_{\operatorname{Out}(\Pi)}
                                      Q_3^{\pm}.             \tag{E.5.2}
\]
Corollaries~4.5 and 4.7 of \cite{CEBoggiCongruence2020}, for
the class of all finite groups, give the faithful pointed
action and its exact Birman sequence.  Its kernel is
$\Pi\cong\operatorname{Inn}(\Pi)$.  Lifting an outer element
and then adjusting its inner automorphism proves surjectivity
onto the displayed pullback, so these results give exactly
(E.5.2).  Adjoining the actual pointed reflection gives the
same identity for the orientation extension.

Lift $b$ in (E.5.2).  Its conjugate of the pointed reflection
is a section belonging to one of the actual target ovals,
by E.1.  Adjusting by an element of $\Pi$ makes these actual
pointed reflections agree.  The faithfulness in (E.5.2) makes
this an equality in the punctured procongruence group.
Applying E.2 with the one real marked point shows that $b$
takes the original oval generator to the target generator
or its inverse.  Repeating for every oval proves
$b\in(Q_3)_A$, the stabilizer of this actual unoriented
multicurve in the procongruence curve complex.

Put $L=\operatorname{Stab}_{\Gamma_3}(A)$, and let $L_0$ fix
each curve and its orientation.  Cutting and replacing the
boundary curves by labelled punctures gives the exact sequence
\[
 1\longrightarrow T\cong\mathbb Z^r\longrightarrow L_0
       \longrightarrow\prod_j\Gamma(S_j)^{\mathrm{pure}}
       \longrightarrow1,                                    \tag{E.5.3}
\]
where $T$ consists of the individual oval Dehn twists.  The
surface types are
\[
\begin{array}{c|c}
 (r,a)&\coprod_j S_j\\\hline
 (1,1)&S_{2,2}\\
 (2,1)&S_{1,4}\\
 (3,1)&S_{0,6}\\
 (2,0)&S_{1,2}\amalg S_{1,2}\\
 (4,0)&S_{0,4}\amalg S_{0,4}.
\end{array}                                                   \tag{E.5.4}
\]
Theorem~4.5 of \cite{CEBoggiProcongruence2014} identifies the
procongruence stabilizer with the closure of $L$ and gives
the exact completed version of (E.5.3), whose kernel is the
injected full $\widehat{\mathbb Z}^{,r}$.  Each component
has genus at most two.  Theorem~1.4(ii) of
\cite{CEBoggiCongruence2020}, with all finite groups as
formation, identifies its procongruence completion with the
full completion; the genus-zero comparison also follows
from Theorem~7.1 of \cite{CEBoggiHyperelliptic2024}.
Completion of (E.5.3) is right exact.  Its left group maps
injectively to the twist kernel just specified, so it is
also left exact.  The comparison of these two exact
sequences is an isomorphism at both ends and therefore in
the middle.  Accounting for the finite permutation and
orientation quotient gives
\[
                     \widehat L\cong(Q_3)_A.                 \tag{E.5.5}
\]
It also proves that $\widehat L\to P_3$ is injective, since
its composite with $q$ is injective.

We verify the torsion theorem hypotheses for this stabilizer,
not for the ambient genus-three group.  Pure sphere mapping
class groups are good by the proof in E.4; the genus-one
one-puncture group is virtually free; the genus-two closed
group is hyperelliptic, so the same finite-index sphere
argument applies.  Adding punctures preserves goodness by
the exact Birman sequence, goodness of surface groups, and
the Hochschild--Serre comparison.  The same proof gives
finite virtual cohomological type.  Hence all the component
groups in (E.5.3) are good, as are their product and the
twist group.  The exact completion already proved for
(E.5.3) gives goodness of $L_0$.  Finite extensions give
goodness and finite virtual two-cohomological type for
$\langle L,\tau\rangle$.

Theorem~A of \cite{CEBoggiZalesskii2025} now identifies the
centralizer of the one element $\tau$ in $\widehat L$ with
the closure of its discrete centralizer.  A discrete
centralizer preserves $A$, so its centralizer in $L$ is
exactly $C$.  Together with containment in $(Q_3)_A$ and
(E.5.5), this proves the quotient centralizer equality.
The copy of $\overline C$ in $P_3$ maps isomorphically onto
it.  The kernel of the full centralizer map is $K_c^\tau$;
the specified copy is a complement, proving (E.5.1).
Thus only the prime-two element case of the cited theorem
has been used.
\end{proof}

\statement{E.6}{Proved}{The free real type has the exact comparison of LM6 C.18}
For the fixed-point-free genus-three involution $\tau$, put
$N_4=S_3/\tau$, $E=\widehat{\pi_1(N_4)}$, and
$\Pi=\widehat{\pi_1(S_3)}$.  Let $\operatorname{Aut}_{Q}(E;\Pi)$
be the continuous automorphisms preserving $\Pi$ whose
restrictions have outer class in $Q_3^{\pm}$.  Then
\[
 \mathcal A:=\operatorname{Aut}_{Q}(E;\Pi)/\operatorname{Inn}(E)
   \cong C_{Q_3^{\pm}}(\bar\tau)/\langle\bar\tau\rangle
   \cong C_{Q_3}(\bar\tau).                                  \tag{E.6.1}
\]
Under this homeomorphism the closure of
$\operatorname{Mod}(N_4)$ equals the closure of
$C_{\Gamma_3}(\tau)$ in $Q_3$.  Equality of this closure
with $\mathcal A$ remains a density question.

\begin{proof}
The orientation-cover inclusion has full induced profinite
topology, as it has index two, so
$1\to\Pi\to E\to C_2\to1$ is exact.  Choose
$e\in E\setminus\Pi$, put $a=\operatorname{Ad}(e)|_\Pi$
and $z=e^2$.  Then $a^2=\operatorname{Inn}(z)$.  If
$b\in\operatorname{Aut}(\Pi)$ has outer class commuting
with $\bar\tau$, there is a unique $d\in\Pi$ with
$bab^{-1}=\operatorname{Inn}(d)a$, since $Z(\Pi)=1$.
Squaring gives $b(z)=d\,a(d)z$.  These two identities
show that $b$ extends to $E$ by $e\mapsto de$; they
preserve conjugation on $\Pi$ and the square of $e$.
The map is continuous and bijective on the two cosets
$\Pi$ and $\Pi e$.  It is unique by the same centre-free
argument.  Thus restriction gives the inverse image of
the outer centralizer.  Quotienting by inner automorphisms
of $\Pi$ removes $\operatorname{Inn}(\Pi)$, and inner
conjugation by $e$ removes the remaining class $\bar\tau$.
Finally
$C_{Q_3^{\pm}}(\bar\tau)=C_{Q_3}(\bar\tau)\times
\langle\bar\tau\rangle$, by the orientation character.
These constructions give (E.6.1); all continuous bijections
are homeomorphisms by compactness.

Every homeomorphism of $N_4$ has exactly one
orientation-preserving lift to $S_3$.  The lift is compatible
with composition and isotopy and identifies
$\operatorname{Mod}(N_4)$ with $C_{\Gamma_3}(\tau)$.
The precise mapping-class comparison, including its kernel
before choosing the positive lift, is recorded on page 1030
of \cite{CEStukow2016}.  Its completed action is exactly
the restriction in (E.6.1).  Taking closures proves the
last assertion without asserting the density itself.
\end{proof}

\statement{E.7}{Proved}{The obstruction to a hyperelliptic reduction in LM6 C.19}
For the positive lift
$j:\operatorname{Mod}(N_4)\cong C_{\Gamma_3}(\tau)$ and
every finite-index subgroup $H\leq\operatorname{Mod}(N_4)$,
the centralizer $C_{\Gamma_3}(j(H))$ has no nonidentity
finite-order element.  In particular $j(H)$ is never
contained in a hyperelliptic centralizer, and the standard
hyperelliptic subgroup of $\operatorname{Mod}(N_4)$ has
infinite index.

\begin{proof}
Glue two pairs of pants $P,Q$ along boundary circles
$a,b,c$, with orientation voltages $0,0,1$ on the three
edges of the resulting theta graph.  The orientation
double is connected, since the $a,c$ cycle has nonzero
voltage.  The base is a closed nonorientable surface of
Euler characteristic $-2$, hence $N_4$, and its double
is $S_3$.  The lifted pants graph has edges
\[
\begin{array}{c|c}
 a_+,b_+&P_+,Q_+\\
 a_-,b_-&P_-,Q_-\\
 c_+&P_+,Q_-\\
 c_-&P_-,Q_+.
\end{array}                                                   \tag{E.7.1}
\]
The six curves form a pants decomposition.  They are
essential and pairwise nonisotopic: a disc or annulus
bounded by them would be a union of complementary pairs
of pants, contradicting its Euler characteristic.

Choose orientations of the three base annuli and denote
their twists by $t_a,t_b,t_c$.  Their positive lifts are
\[
 j(t_a)=T_{a_+}T_{a_-}^{-1},\quad
 j(t_b)=T_{b_+}T_{b_-}^{-1},\quad
 j(t_c)=T_{c_+}T_{c_-}^{-1}.                                  \tag{E.7.2}
\]
The opposite signs arise because the deck involution
exchanges the two annuli and reverses the surface
orientation.  A finite-index $H$ contains positive powers
of all three twists.  If a positive finite-order mapping
class $f$ centralizes their lifts, uniqueness of signed
multitwist supports and exponents forces it to fix each
of the six individual curves.  The two exponents $m$ and
$-m$ cannot be interchanged by a positive mapping class.
This is precisely Lemma~3.17 of \cite{CEFarbMargalit2011}.

Each vertex in (E.7.1) has a distinct set of incident
labelled edges, so $f$ also fixes every complementary
pair of pants.  Its restrictions, with boundaries treated
as punctures, lie in the trivial pure three-punctured-sphere
mapping class groups.  By the cutting homomorphism,
Proposition~3.20 of the same source, $f$ lies in the
multitwist kernel $\mathbb Z^6$.  This group is torsion-free,
so $f=1$.  Proposition~2.3 there supplies the asserted
triviality for the pure three-punctured sphere.  The
standard nonorientable hyperelliptic subgroup lifts into
a hyperelliptic centralizer by the explicit maps on page
1030 of \cite{CEStukow2016}; it therefore has infinite
index.  The graph and the signed exponent vectors are
independently checked by \texttt{code/check\_real\_models.py}.
\end{proof}

\statement{E.8}{Proved}{Finite quotients retain the entire fixed group and nonabelian cohomology}
Let a finite group $C_2=\langle\tau\rangle$ act continuously
on a profinite group $K$, and let $N$ range over its open
normal $\tau$-invariant subgroups.  There are natural
homeomorphisms
\[
 K^\tau\cong\varprojlim_N(K/N)^\tau,
 \qquad
 H^1(C_2,K)\cong\varprojlim_N H^1(C_2,K/N).                  \tag{E.8.1}
\]
The cohomology set is given the quotient topology of its
compact cocycle space.  Neither the maps to individual
finite fixed groups nor the finite cohomology transition
maps are asserted to be surjective.

\begin{proof}
The fixed-group identity follows by applying the defining
equation $\tau(k)=k$ in every finite quotient.  For
cohomology, a compatible family of finite cohomology
classes defines, for each $N$, a nonempty closed subset
of $K$ consisting of all lifts of its cocycle class in
$K/N$.  Any finitely many such conditions are satisfied
by a lift from one common refinement, so compactness
gives $k\in K$ satisfying them all.  Its cocycle equation
holds modulo every $N$ and hence in $K$ itself.  This
proves surjectivity.  If two actual cocycles have the same
finite class at every $N$, the sets of conjugators
$b\in K$ satisfying $k'=bk\tau(b)^{-1}\pmod N$ are
nonempty closed sets with the same finite-intersection
property.  Their intersection gives an actual
conjugator, proving injectivity.  The orbit relation
is closed, as it is the image of a compact set, so
the cohomology space is compact Hausdorff.  The continuous
bijection to the inverse limit is therefore a
homeomorphism.
\end{proof}

\statement{E.9}{Proved}{The fixed congruence kernel is the fibre of oval stabilizers}
For one of the five oval types in E.5, let $A$ be its oval
multicurve and $L=\operatorname{Stab}_{\Gamma_3}(A)$.
Use E.5 to regard $\widehat L$ as a closed subgroup of
$P_3$ mapping isomorphically onto $(Q_3)_A$.  Then the
fibre over the identity coset of the continuous map
\[
 P_3/\widehat L\longrightarrow Q_3/(Q_3)_A
\]
is naturally the profinite set $K_c$.  Its $\tau$-fixed
subset is naturally $K_c^\tau$.

\begin{proof}
The inverse image of $(Q_3)_A$ is
$K_c\widehat L$, and $K_c\cap\widehat L=1$.  Every coset
in this fibre therefore has a unique representative
$k\in K_c$.  This is a homeomorphism by compactness.
The actual involution $\tau$ normalizes $L$ and its
completion.  Its action on the fibre takes $k\widehat L$
to $\tau(k)\widehat L$.  This equals $k\widehat L$
exactly when $k^{-1}\tau(k)\in K_c\cap\widehat L=1$.
This proves the fixed-subset assertion.  In particular,
a theorem saying that there is just one fixed lift of
this geometric multicurve would prove $K_c^\tau=1$;
preservation of the procongruence multicurve alone does
not say that.
\end{proof}

\statement{E.10}{Open}{The six congruence-kernel computations in prompt E.1}
For each type (E.1.1), determine
$K_c^\tau$ and the nonabelian pointed set
$H^1(C_2,K_c)$ for $\operatorname{Ad}(\tau)$.
E.3 gives the exact class and centralizer formula after
these groups and sets are known; E.5 supplies its
quotient-centralizer term for the five oval types.
For $(0,1)$ the density question in E.6 remains an
additional unresolved term.  The current state is
\[
\begin{array}{c|c|c|c}
 (r,a)&C_{Q_3}(\bar\tau)&K_c^\tau&H^1(C_2,K_c)\\\hline
 (0,1)&\mathcal A\text{ of E.6}&\text{uncomputed}&\text{uncomputed}\\
 (1,1)&\overline C^{\,Q_3}&\text{uncomputed}&\text{uncomputed}\\
 (2,1)&\overline C^{\,Q_3}&\text{uncomputed}&\text{uncomputed}\\
 (3,1)&\overline C^{\,Q_3}&\text{uncomputed}&\text{uncomputed}\\
 (2,0)&\overline C^{\,Q_3}&\text{uncomputed}&\text{uncomputed}\\
 (4,0)&\overline C^{\,Q_3}&\text{uncomputed}&\text{uncomputed}.
\end{array}                                                   \tag{E.10.1}
\]
The existence of six geometric classes is proved, and
their distinctness persists in both completions; a
classification of all profinite classes does not follow.
E.8 gives exact inverse-limit computations but no
presentation of the required finite quotients of $K_c$.
E.9 identifies a separate geometric fixed-fibre problem.
The hyperelliptic route settles the intersection in
E.4; E.7 rules out extending that route to the free
type by passing to a finite-index nonorientable
mapping class subgroup.  No vanishing of the genus-three
congruence kernel is assumed.  The contemporary source
\cite[Problem~19 and the following paragraph]{CEAndrewHenselHughesWade2026}
explicitly records that the mapping class congruence
subgroup problem remains open in genus at least three.

